\documentclass[main.tex]{subfiles}
\newcommand{\DiagramaEscala}{0.9}
\begin{document}

\section{Ecuaciones Diferenciales Parciales}

% --- Diapositiva 1: Definición de derivada parcial ---
\begin{frame}{Motivación: Derivadas Parciales}
\small
Dada una función \( u(x, y) \), la derivada parcial con respecto a \( x \) se define como:
\[
\frac{\partial u}{\partial x} = \lim_{\Delta x \to 0} \frac{u(x + \Delta x, y) - u(x, y)}{\Delta x}
\]
De forma similar, la derivada parcial con respecto a \( y \) es:
\[
\frac{\partial u}{\partial y} = \lim_{\Delta y \to 0} \frac{u(x, y + \Delta y) - u(x, y)}{\Delta y}
\]
\end{frame}

% --- Diapositiva 2: ¿Qué es una EDP? ---
\begin{frame}{¿Qué es una Ecuación Diferencial Parcial (EDP)?}
\small
Una EDP es una ecuación que contiene derivadas parciales de una función desconocida de dos o más variables independientes.

\textbf{Ejemplos:}
\begin{align*}
\frac{\partial^2 u}{\partial x^2} + x y \frac{\partial^2 u}{\partial y^2} + u &= 1 \\
\frac{\partial^3 u}{\partial x^2 \partial y} + 8x \frac{\partial u}{\partial y} + 5u &= y \\
\frac{\partial^2 u}{\partial x^2} + \frac{\partial^3 u}{\partial x \partial y^2} &= 6 \\
\frac{\partial^2 u}{\partial x^2} + x \frac{\partial u}{\partial x} &= \frac{\partial u}{\partial y}
\end{align*}
\end{frame}

% --- Diapositiva 3: Orden y linealidad de una EDP ---
\begin{frame}{Orden y Linealidad}
\small
\textbf{Orden:} Determinado por la derivada de mayor orden presente.

\textbf{Ejemplos:}
\begin{itemize}
  \item PT8.3: Orden 2
  \item PT8.4: Orden 3
\end{itemize}

\textbf{Linealidad:} Una EDP es lineal si es lineal en \( u \) y sus derivadas, con coeficientes que dependen solo de las variables independientes.

\textbf{Ejemplos:}
\begin{itemize}
  \item PT8.3 y PT8.4: lineales
  \item PT8.5 y PT8.6: no lineales
\end{itemize}
\end{frame}

% --- Clasificación de EDPs ---
\begin{frame}{Clasificación de EDPs de segundo orden}
\small
Forma general:
\[
A \frac{\partial^2 u}{\partial x^2} + B \frac{\partial^2 u}{\partial x \partial y} + C \frac{\partial^2 u}{\partial y^2} + D = 0
\]
\textbf{Clasificación según } \( B^2 - 4AC \):
\begin{itemize}
  \item \( < 0 \): Elíptica
  \item \( = 0 \): Parabólica
  \item \( > 0 \): Hiperbólica
\end{itemize}
\end{frame}

% --- Ejemplos de cada categoría ---
\begin{frame}{Ejemplos de Clasificación}
\small
\begin{itemize}
  \item Elíptica: \( \frac{\partial^2 T}{\partial x^2} + \frac{\partial^2 T}{\partial y^2} = 0 \)
  \item Parabólica: \( \frac{\partial T}{\partial t} = k' \frac{\partial^2 T}{\partial x^2} \)
  \item Hiperbólica: \( \frac{\partial^2 y}{\partial x^2} = \frac{1}{c^2} \frac{\partial^2 y}{\partial t^2} \)
\end{itemize}
\end{frame}

% --- Motivación física: placa caliente ---
\begin{frame}{29.1 La ecuación de Laplace: Motivación Física}
\small
Placa rectangular delgada con espesor \(\Delta z\), aislada excepto en los extremos.

Flujo neto de calor en estado estacionario:
\[
q(x) \Delta y \Delta z \Delta t + q(y) \Delta x \Delta z \Delta t = q(x + \Delta x) \Delta y \Delta z \Delta t + q(y + \Delta y) \Delta x \Delta z \Delta t
\]
Dividiendo entre \( \Delta z \Delta t \) y reorganizando:
\[
[q(x) - q(x + \Delta x)] \Delta y + [q(y) - q(y + \Delta y)] \Delta x = 0
\]
\end{frame}

% --- De flujo a derivadas ---
\begin{frame}{Derivadas parciales del flujo}
\small
Reescribiendo con cocientes de diferencia:
\[
\frac{q(x) - q(x + \Delta x)}{\Delta x} + \frac{q(y) - q(y + \Delta y)}{\Delta y} = 0
\]
Tomando el límite cuando \( \Delta x, \Delta y \to 0 \):
\[
-\frac{\partial q}{\partial x} - \frac{\partial q}{\partial y} = 0
\]
\end{frame}

% --- Ley de Fourier ---
\begin{frame}{Ley de Fourier del calor}
\small
\textbf{Ley de Fourier:} El flujo de calor es proporcional al gradiente de temperatura:
\[
q_i = -k \rho C \frac{\partial T}{\partial x_i}
\]
\textbf{Definición de conductividad térmica:}
\[
k' = k \rho C
\]
\textbf{Sustituyendo en la ecuación:}
\[
\frac{\partial^2 T}{\partial x^2} + \frac{\partial^2 T}{\partial y^2} = 0 \quad \text{(Ecuación de Laplace)}
\]
\end{frame}

% --- Ecuación de Poisson ---
\begin{frame}{Ecuación de Poisson}
\small
Si existen fuentes o pérdidas de calor:
\[
\frac{\partial^2 T}{\partial x^2} + \frac{\partial^2 T}{\partial y^2} = f(x, y)
\]
\textbf{f(x,y)} representa las fuentes internas (positiva) o sumideros (negativa) de calor en el dominio.

\textbf{Ecuaciones:}
\begin{itemize}
  \item \textbf{Laplace:} sin fuentes internas
  \item \textbf{Poisson:} con fuentes internas
\end{itemize}
\end{frame}


\end{document}